\documentclass[10pt,a4paper]{amsart}
\usepackage[margin=19mm]{geometry}
\usepackage{amsmath,amssymb,mathtools}
\usepackage[scr=rsfs,cal=euler]{mathalfa}
\usepackage{xfrac}
\usepackage[unicode,hidelinks,bookmarks=false]{hyperref}
\newcommand{\ie}{i.e.\ }
\newcommand{\cf}{cf.\ }
\newcommand{\resp}{respectively\ }
\newcommand{\Subsection}{\subsection}
\providecommand{\nobreakdash}{\nobreak\hbox{-}}
\setlength{\emergencystretch}{2em}
\multlinegap0pt
\usepackage{amssymb}
\usepackage[shortlabels]{enumitem}
\usepackage[all]{xy}
\usepackage{float}
\usepackage{xcolor}
\newtheorem{theorem}{Theorem}
\title{The steady state of gravity-capillary problem with inclined walls}
\author{Xiaoding Yang}
\date{Source: arXiv:2503.04048v2 (CC BY 4.0)}
\begin{document}
\maketitle

\noindent\emph{Source and licence.} The steady state of gravity-capillary problem with inclined walls. Version arXiv:2503.04048v2 (CC BY 4.0), distributed by arXiv under the Creative Commons Attribution 4.0 International licence, http://creativecommons.org/licenses/by/4.0/, as recorded in the arXiv licence metadata for this exact version and shown on the abstract page https://arxiv.org/abs/2503.04048v2. Full licence notice: \texttt{licenses/arxiv-2503.04048-CC-BY-4.0.txt}.\par\smallskip

\noindent\textit{Editorial scope.} Counted unit: the article's theorem theorem (source0.tex lines 1689--1705). The article's complete original proof of it is delivered with it (source0.tex lines 1707--1893, one \texttt{proof} environment of 187 lines). No mathematics is written, repaired or re-derived: the statement and the proof are the article's own lines, in the article's own notation, and no retained line is substituted or re-flowed.  Retained uncounted as the source-local context the proof needs: lines 209--216; lines 1642--1644.  Nothing was omitted from any retained span, so the package carries no omission record.  No retained line contains a citation, so the wrapper carries no bibliography and no bibliography entry.  No retained line contains a figure environment, an image-inclusion command or drawing code, so no image is delivered and the package declares no asset dependency.  Editorial formatting only, in addition: an AMS \texttt{amsart} preamble that restates the article's own declarations and macros used by the retained lines, namely the environments \texttt{theorem} on separate counters, together with the standard packages those lines need.  The retained environments print in this document as Theorem 1 in order of appearance, whereas the article prints the same items with the numbering of its own full document.  The complete native source of the article as distributed in the arXiv e-print for this exact version is bundled unchanged as \texttt{sources/174301/source0.tex}.
     \begin{equation}{\label{equ:1.2.600}}
    \begin{cases}
        \frac{dx}{d\psi}=\sigma\frac{\cos\psi}{gy-P_{0}}\\
    \frac{dy}{d\psi}=\frac{dy}{dx}\frac{dx}{d\psi}=\sigma\tan\psi\frac{\cos\psi}{gy-P_{0}}=\sigma\frac{\sin\psi}{gy-P_{0}}\\
    \sin\psi(\pi-\theta_1)=\sin\psi_1=(-\frac{[\![\gamma]\!]}{\sigma})\sin\theta_1+\sqrt{1-\frac{[\![\gamma]\!]^{2}}{\sigma^{2}}}\cos\theta_1\\
    \sin\psi(\theta_2)=\sin\psi_{2}=(\frac{[\![\gamma]\!]}{\sigma})\sin\theta_2+\sqrt{1-\frac{[\![\gamma]\!]^{2}}{\sigma^{2}}}\cos\theta_2
    \end{cases}
    \end{equation}

     \begin{equation}{\label{equ:3.1.1}}
         gv(x)-\sigma \mathcal{H}=P_{0},
     \end{equation}

    \begin{theorem}{The shape of the function }

Assume the same hypotheses as in Theorem 3.1. Then the following statements hold:
    
    1: If
    \begin{align}
        gV+\sigma \sin(\gamma)-\sigma \sin(-\frac{\pi}{2}-\gamma)\leq 0, 
    \end{align}
     then $\psi$ is a monotone increasing function of $x$.

    2:If

    \begin{align}
        gV+\sigma \sin(\gamma)-\sigma \sin(-\frac{\pi}{2}-\gamma)>0,
    \end{align}
      then $\psi$ attains its unique maximal value $\psi_m$ in the interior of the interval where $ \psi_m\in(\max(\gamma-\frac{\pi}{2},-\gamma),0)$. 
    \end{theorem}

    \begin{proof}

    \textbf{proof of} (1):  Integrating both sides of the equation \eqref{equ:3.1.1} from 0 to $x_2$, we obtain

    \begin{equation}{\label{equ:special}}
        gV+\sigma \sin(\frac{\pi}{2}-\gamma)-\sigma\sin(\gamma)=x_2P_{0}
    \end{equation}

    \noindent Using the condition $gV+\sigma \sin(\phi)-\sigma \sin(\frac{\pi}{2}-\phi)\leq 0$, we derive $P_{0}\leq 0$ from equation \eqref{equ:special}. Substituting this into \eqref{equ:3.1.1} gives
    
    \begin{equation*}
        \sigma \partial_{x} \sin\psi=gv(x)-P_{0}>0~\operatorname{for~any}~x.
    \end{equation*}

    \noindent Hence, $\psi$ is an increasing function of $x$, and therefore the curve can be parameterized by $\psi$.

    \begin{equation*}
        ~
    \end{equation*}
    
    \textbf{proof of} (2): From Theorem 3.1, $P_{0}> 0$ when  $gV+\sigma \sin(\phi)-\sigma \sin(\frac{\pi}{2}-\phi)>0$. In this case $\psi$ is neither monotonically decreasing nor increasing. We divide the main proof into several steps.

    \textbf{Step 1} In this step, we prove that the boundary curve in this case must attain its maximum $\psi_{m}$ at some point. Moreover, $\psi_m\leq0$. 
    
    To prove that the boundary curve attains its maximum at some point. It suffices to show the following property

    \begin{equation}{\label{key_assume}}
        \sigma \frac{d}{dx}\sin(\psi)|_{x=0}=gv(x_1)-P_{0}>0.
    \end{equation}

    \noindent We argue by contradiction and assume instead that

    \begin{equation}{\label{equ:3.1.3}}
        \sigma \frac{d}{dx}\sin(\psi)|_{x=0}=gv(x_1)-P_{0}\leq 0
    \end{equation}

     Differentiating both sides of equation \eqref{equ:3.1.1} with respect to $x$, we obtain

    \begin{equation}{\label{equ:3.1.4}}
        \sigma \frac{d^{2}}{dx^{2}}\sin(\psi)|_{x=0}=g\partial_{x}v|_{x=0}=g \tan\psi|_{x=0}=g \tan (\gamma)< 0.
    \end{equation}

    \noindent Consequently,

    \begin{equation*}
        \sigma \frac{d^{2}}{d^{2}x}\sin(\psi)<0,
    \end{equation*}

    \noindent in a neighborhood of $x=0$, which implies

    \begin{equation*}
        \sin\psi(x)<0.
    \end{equation*}
    
    \noindent for $x$ close enough to $0$. We now define

    \begin{equation*}
        x_{0}:=\min\{x:\frac{d^{2}}{dx^{2}}\sin\psi=0\}.
    \end{equation*}

    \noindent This point is well-defined since, by assumption, $\psi$ is not a monotone decreasing function of $x$. From the computation in equation \eqref{equ:3.1.4}, we have $\tan\psi=0$ when $x=x_0$. Hence, $\sin\psi(x_0)=0>\sin\psi_1$. Combining this result with the assumption \eqref{equ:3.1.3}, we deduce that there exits a point $x_{1}\in(x,x_{0})$ such that $\frac{d}{dx}(\sin\psi)|_{x=x_{1}}=0$. However, by the definition of $x_{0}$, we know that 
    \begin{align}
        \frac{d^{2}}{dx^{2}}(\sin\psi)<0
    \end{align}
    for all $0<x<x_0$. Therefore, $\frac{d}{dx}(\sin\psi)<0$ for all $x\in(0,x_0)$. This  contradicts the existence of such points $x_0$ and $x_{1}$. Therefore, assumption \eqref{equ:3.1.3} can not hold.
    
    We now show that the maximum angle $\psi_{m}\leq0$. From \eqref{key_assume}, we obtain the following property

    \begin{equation}{\label{equ:3.1.5}}
        \sigma \frac{d}{d\psi}\sin\psi|_{x=x_{1}}=gv(x_{1})-P_{0}>0
    \end{equation}

    \noindent Suppose, for the sake of contradiction, that the maximum value $\psi_m>0$. Then we have:

    \begin{equation*}
        \sigma \frac{d}{d\psi}\sin\psi|_{x=x_m}=gv(x_m)-P_{0}=0
    \end{equation*}

    \noindent From equation \eqref{equ:3.1.4}, we have

    \begin{equation}{\label{equ:3.1.6}}
        \sigma \frac{d^{2}}{d\psi^{2}}\sin(\psi)|_{x=x_m}=g\tan\psi_m>0
    \end{equation}

    \noindent However, since $\psi_m$ is assumed to be a maximum, we must have
    
    \begin{equation*}
        \sigma \frac{d^{2}}{d\psi^{2}}\sin(\psi)|_{\psi=\psi_m}=g\tan\psi_m\leq 0,
    \end{equation*}

    \noindent which contradicts \eqref{equ:3.1.6}. Therefore, the maximum for the inclined angle $\psi_{m}$ must satisfy $\psi_m\leq 0$.
    
     \textbf{Step 2}: In this step, we prove that $\psi_m\neq 0$. 
     
     We proceed by a contradiction argument again. Suppose that $\psi_m=0$. From the first and second equation of the system \eqref{equ:1.2.600}, we have

    \begin{equation}{\label{equ:3.1.7}}
        \frac{d\psi}{dx}=\frac{gv-P_{0}}{\sigma \cos\psi}
    \end{equation}

    \begin{equation}{\label{equ:3.1.8}}
        \frac{dy}{d\psi}=\frac{\sigma \sin\psi}{gv-P_{0}}
    \end{equation}

    \noindent Integrating both sides of the equation \eqref{equ:3.1.8} from $\psi$ to $\psi_m$, we obtain

    \begin{equation}{\label{equ:3.1.9}}
        y(\psi)=\frac{P_{0}\pm \sqrt{P_{0}^{2}-2g\sigma(\cos\psi-C)}}{g},
    \end{equation}

    \noindent where $C$ is a constant determined by condition $v(\psi_m)=\frac{P_{0}}{g}$.

    Plugging the equation \eqref{equ:3.1.9} into the ODE \eqref{equ:3.1.7}, we obtain

    \begin{equation}{\label{equ:3.1.110}}
         F(\psi):=\frac{d\psi}{dx}=\frac{\sqrt{P_{0}^{2}-2g\sigma(\cos\psi-C)}}{\sigma \cos\psi}.
    \end{equation}
    
    \noindent Differentiating function $F$ with respect to $\psi$ near $\psi_m$,  we obtain

    \begin{equation}{\label{equ:3.1.10}}
        \frac{dF(\psi)}{d\psi}=-\frac{\sin\psi\sqrt{P_{0}^{2}-2g\sigma(\cos \psi-C)}}{\sigma \cos^{2}\psi}+\frac{d\sqrt{P_{0}^{2}-2g\sigma(\cos \psi-C)}}{d\psi}\frac{1}{\sigma \cos \psi}.
    \end{equation}

    \noindent  The first term of equation \eqref{equ:3.1.10} is a smooth function that tends to zero as $\psi\rightarrow 0$, and is therefore bounded in a neighborhood of $\psi_m=0$. For the second term, we have
    
    \begin{equation}{\label{equ:3.1.11}}
        \frac{d\sqrt{P_{0}^{2}-2g\sigma(cos \psi-C)}}{d\psi}=-\frac{1}{2}\frac{2g\sigma \sin\psi}{\sqrt{P_{0}^{2}-2g\sigma(\cos \psi-C)}}.
    \end{equation}

     \noindent Performing a Taylor expansion of the right-hand side of equation \eqref{equ:3.1.11} near $\psi=0$, we obtain

    \begin{align}
        \frac{d\sqrt{P_{0}^{2}-2g\sigma(\cos \psi-C)}}{d\psi}&=-\frac{1}{2}\frac{2g\sigma \sin\psi}{\sqrt{P_{0}^{2}-2g\sigma(\cos \psi-C)}} \notag\\
        &=-g\sigma \frac{\psi+O(\psi^{2})}{\sqrt{2g\sigma}\sqrt{\frac{P_{0}^{2}}{2g\sigma}+C-\cos\psi}} \label{equ:3.1.12}.
    \end{align}

     From the definition of constant $C$, we have

    \begin{equation*}
        \frac{P_{0}^{2}}{2g\sigma}+C-\cos\psi_m=0
    \end{equation*}

    \noindent Hence, the denominator in \eqref{equ:3.1.12} can be expanded as

    \begin{equation*}
        \sqrt{\frac{P_{0}^{2}}{2g\sigma}+C-\cos\psi}=\sqrt{1-\cos\psi}=\frac{\sqrt{2}}{2}\psi+O(\psi^{2})
    \end{equation*}
    
    \noindent  Substituting it into the derivative term \eqref{equ:3.1.12}, we obtain

    \begin{align}
         \eqref{equ:3.1.12}&=-g\sigma \frac{\psi+O(\psi^{2})}{\sqrt{2g\psi}\sqrt{\frac{P_{0}^{2}+C}{2g\psi}-\cos\psi}} \notag\\
         &= -g\sigma \frac{\psi+O(\psi^{2})}{\sqrt{2g\sigma}\sqrt{1-\cos\psi}} \notag\\
         &=-g\sigma \frac{\psi+O(\psi^{2})}{\sqrt{g\sigma}\psi+O(\psi^{2})} \notag\\
          &=-\sqrt{g\sigma} (1+O(\psi)) \label{equ:3.1.13}
    \end{align}
    
      Thus, for the ODE \eqref{equ:3.1.110} written as follows

    \begin{equation*}
        \frac{d\psi}{dy}=\frac{\sqrt{P_{0}^{2}-2g\sigma(\cos\psi-C)}}{\sigma \cos\psi}=F(\psi),
    \end{equation*}

    \noindent we infer from \eqref{equ:3.1.10} and \eqref{equ:3.1.13} that $F(\psi)$ has a bounded first derivative and therefore Lipschitz continuous in a neighborhood of $\psi=\psi_m=0$. By the uniqueness theorem for ODEs,  $\psi=0$ then must be the unique solution of the ODE, which contradicts our assumption that $\psi_{m}$ is a nontrivial maximum.

    \textbf{Step 3} We prove the uniqueness of the maximum in this step. 
    
    From Step 1 and Step 2, $\psi$ reaches to a maximum $\psi_{m}<0$. Since
    \begin{align}
        \frac{d}{dx}\sin\psi=gv(x)-P_{0},
    \end{align}
    we have
    \begin{align}{\label{equ:negative}}
        \frac{d}{dx}\sin\psi|_{x=x_{m}}=gv(x_{m})-P_{0}=0.
    \end{align}
    \noindent Since $\psi_{m}<0$, the function $v(x)$ is monotonically decreasing in a neighborhood of $x_{3}$. Therefore, we have
    \begin{align}
        \frac{d}{dx}\sin\psi=gv(x)-P_{0}<0
    \end{align}
    for all $x\in (x_{m},x_{m}+\delta)$ for some small $\delta$. We can extend this argument inductively to conclude that the relation in \eqref{equ:negative} is true for all $x\in (x_{m},x_{4})$ where $x_{4}$ is a point such that $x_{4}>x_{m}$ and $\psi(x_{4})=\psi_{2}$. 
    
    If the curve $(x,v(x))$ does not touch the right wall at $x= x_{4}$, then there exists a point $x_{5}>x_{4}$ such that $\psi(x_{5})=-\frac{\pi}{2}$ which contradicts the assumption that the curve can be expressed as a function of $x$. Therefore, the point $(x_{4},y(x_{4}))=(x_{2},u_{2})$ is the contact point with the right wall. 

    From the discussion above, $\psi(x)$ is monotonically increasing when $x\in (x_{1},x_{m})$, and then monotonically decreasing when $x\in (x_{m},x_{2})$. This implies that the maximum $\psi_{m}=\psi(x_{m})$ is unique.
    
    \end{proof}

\end{document}
